%! AP Statistics — Unit 2, Lesson 2.9 — Exit Ticket — Answer Key
\documentclass[10pt]{article}
\usepackage{apstats-article}
\usepackage{apstats-key}

\begin{document}

\pageheader{Unit 2, Lesson 2.9}{Exit Ticket --- Answer Key}
\namedateperiod

\vspace{0.18in}

% ── SCENARIO ─────────────────────────────────────────────────────────────────
\begin{tcolorbox}[colback=sky, colframe=navy, boxrule=0.9pt, arc=2mm,
    left=4mm, right=4mm, top=2.5mm, bottom=2.5mm]
\small\textit{A marine biologist recorded the age of blue whales (years, $x$) and their body
mass (tonnes, $y$). The scatterplot curved upward. After a $\log_{10}$ transformation of $y$,
the new scatterplot was linear. Technology output:
$\;\widehat{\log(\text{mass})} = 0.95 + 0.065x$, with $r^2 = 0.91$.}
\end{tcolorbox}

\vspace{0.16in}

\begin{enumerate}[leftmargin=1.6em, itemsep=0.16in]

  \item State which variable was transformed and what type of model (exponential or power)
        this implies.\\[8pt]
        \ansline{Only $y$ (mass) was transformed. Transforming only $y$ produces an
        \textbf{exponential model}: $\text{mass} = a \cdot b^{\text{age}}$.}

  \item Write the back-transformed equation for predicting whale mass on the original scale.
        Show your work.\\[8pt]
        \ansline{$\widehat{\text{mass}} = 10^{0.95 + 0.065x}$\enspace
        (equivalently $\approx 8.91 \cdot 1.162^x$ tonnes)}

  \item Predict the mass of a blue whale that is \textbf{10 years old}. Show all work,
        including the back-transformation step.\\[8pt]
        \ansline{$\widehat{\log(\text{mass})} = 0.95 + 0.065(10) = 0.95 + 0.65 = 1.60$}
        \ansline{$\widehat{\text{mass}} = 10^{1.60} \approx 39.8$ tonnes}

\end{enumerate}

\vspace{0.12in}

% ── AP-STYLE MC ──────────────────────────────────────────────────────────────
\begin{tcolorbox}[colback=redbg, colframe=redacc, boxrule=0.8pt, arc=2mm,
    left=4mm, right=4mm, top=2.5mm, bottom=2.5mm]
\small\textbf{AP-Style Practice}\enspace
A student fits a regression of $\log y$ on $x$ and gets $\widehat{\log y} = 1.2 + 0.15x$.
To predict $y$ when $x = 4$, which calculation is correct?

\vspace{7pt}
\begin{enumerate}[label=\textbf{(\Alph*)}, leftmargin=2em, itemsep=4pt]
  \item $\hat{y} = 1.2 + 0.15(4) = 1.80$
  \item $\hat{y} = 10^{1.2} + 0.15(4) = 16.55$
  \item $\hat{y} = 10^{1.2 + 0.15(4)} = 10^{1.80} \approx 63.1$
        \textcolor{keyred}{\textbf{$\leftarrow$ correct}}
  \item $\hat{y} = \ln(1.2 + 0.15 \times 4) = \ln(1.80) \approx 0.588$
\end{enumerate}
\end{tcolorbox}

\vspace{0.1in}
\begin{itemize}[leftmargin=1.4em]
  \item \textbf{(A)} is wrong: reports the transformed value $\widehat{\log y}$, not
        back-transformed $\hat{y}$.
  \item \textbf{(B)} is wrong: incorrectly mixes back-transformation and the linear term.
  \item \textbf{(C)} is correct: compute $\widehat{\log y} = 1.80$, then back-transform
        $\hat{y} = 10^{1.80} \approx 63.1$.
  \item \textbf{(D)} is wrong: uses natural log instead of base-10 log.
\end{itemize}

\end{document}
